\documentclass[10pt]{article}
\usepackage[T1]{fontenc}
\usepackage{lmodern}
\usepackage{amsmath,amssymb,amsthm,mathtools}
\usepackage{geometry,booktabs,array}
\usepackage{microtype}
\usepackage[hidelinks]{hyperref}
\usepackage{fancyhdr}
\geometry{margin=0.80in,headheight=14pt}
\pagestyle{fancy}
\fancyhf{}
\fancyhead[L]{\small No.72 / Negative Research Record}
\fancyhead[R]{\small v1.0 / 18 Nov 2026}
\fancyfoot[C]{\thepage}
\newtheorem{theorem}{Theorem}
\newtheorem{proposition}[theorem]{Proposition}
\newtheorem{corollary}[theorem]{Corollary}
\newtheorem{lemma}[theorem]{Lemma}
\newcommand{\Ga}{\Gamma}
\newcommand{\capt}{\operatorname{capt}}
\newcommand{\dist}{\operatorname{dist}}
\title{A Failed Static-Certificate Route to Capture-Time Bounds on Fibonacci Cubes\\[3pt]
\large A Qualified Negative Research Record and an Audit}
\author{Ryutaro Yonezu\\\small Independent research / No.72}
\date{18 November 2026 --- v1.0 negative research record}
\begin{document}
\maketitle
\begin{center}\fbox{\parbox{0.86\linewidth}{\centering\small
\textbf{NEGATIVE RESEARCH RECORD -- NOT PEER REVIEWED -- NOVELTY INCONCLUSIVE}\\
Novelty: inconclusive. No claim of a new capture-time formula.}}
\end{center}
\begin{abstract}
We document a failed attempt to infer short capture-time lower bounds on Fibonacci
cubes from static distance covering and a delayed-retract condition.
An elementary construction covers every Fibonacci cube $\Ga_{2m}$ by three
balls of radius $m-1$, and every $\Ga_{2m+1}$ ($m\ge2$) by four balls
of the same radius. Together with the established exact cop-number formula, these yield
infinite families of finite-radius coverings by fewer centers than the number
of cops required for eventual capture. This comparison is elementary and
\emph{does not show} that the displayed radii are below the minimum-cop
capture times in those families; even a single center has finite covering
radius. At dimension ten, a three-center radius-four covering coexists with the independently
audited computer-assisted value $\capt(\Ga_{10})=5$ for three cops.
At dimension eleven, four explicit centers satisfy both the radius-four
covering requirement and the particular $Q_6$ delayed-retract test, without
settling whether four-move capture is possible.
The result is a precise limitation on \emph{these necessary conditions}, not an
impossibility theorem for all approaches, and not a proof that
$\capt(\Ga_{11})\ge5$.
\end{abstract}

\section{Question and prior art}
Let $\Ga_n$ denote the induced subgraph of the $n$-dimensional hypercube
on binary strings containing no adjacent $1$'s. We identify a string with
the integer whose bit $i$ is its coefficient of $2^i$, with $i=0$ the
least significant coordinate. Cops choose their positions before the robber;
subsequently the cops move first. Every player may move along one edge or stay.
We write $c(G)$ for the minimum number of cops with an eventual winning strategy,
and $\capt_k(G)$ for the least guaranteed number of police moves when $k$ cops
are available. The usual $\capt(G)$ is $\capt_{c(G)}(G)$.

Huang, Xie and Xu proved the exact formula \cite{hxx2026}
\begin{equation}\label{eq:cop}
 c(\Ga_n)=\left\lceil\frac{\lceil n/2\rceil+1}{2}\right\rceil.
\end{equation}
The capture-time asymptotic for hypercubes is established by
Bonato, Gordinowicz, Kinnersley and Pra\l at \cite{bgkp2013}.
Retracts and finite-horizon capture are established tools in this area
\cite{bonato2009,over2016}. We therefore do \emph{not} claim that using
distance covering, a retraction, or a hitting-set certificate is a new technique.

Research No.72 had previously established computer-assisted small-dimensional
results, including $\capt_3(\Ga_9)=4$ and $\capt_3(\Ga_{10})=5$,
with independent reimplementations and a separately audited Windows evidence
package (v0.18). Its earlier purported five-move lower bound for $\Ga_{11}$
was withdrawn after the last 616 negative instances were found invalid:
a legal empty responsibility choice for a cop had been omitted. The present
note \emph{does not revive} that lower bound.

Related structural studies determine the radius and center of Fibonacci cubes \cite{munarini2002}; domination-type invariants and covering-code considerations have also been investigated \cite{azarija2018}. These works constrain claims of priority for elementary covering constructions.\par\smallskip

\section{An explicit geometric covering in every dimension}
For $C\subseteq V(G)$, define its static covering radius by
\[
 R_G(C)=\max_{v\in V(G)}\min_{x\in C}d_G(v,x).
\]
The necessary condition for an opening $C$ to win within $t$ police
moves is $R_G(C)\le t$: a stationary robber cannot be reached otherwise.
This necessary condition says nothing about a robber who moves.

\begin{lemma}\label{lem:metric}
For $x,y\in\Ga_n$, the graph distance is the Hamming distance
$d_{\Ga_n}(x,y)=\mathrm{wt}(x\oplus y)$.
\end{lemma}
\begin{proof}
Every edge flips one coordinate, so Hamming distance is a lower bound.
Starting at $x$, first remove the $1$ bits absent from $y$, then add the
$1$ bits present in $y$ and absent from $x$. Deletions preserve legality,
and each intermediate set of added bits is a subset of the legal $1$ set
of $y$. This gives a path of exactly $\mathrm{wt}(x\oplus y)$ edges.
\end{proof}

\begin{theorem}[constant-size ball covers]\label{thm:cover}
For every integer $m\ge1$,
\[
 R_{\Ga_{2m}}(\{0,1,2\})\le m-1.
\]
For every $m\ge2$,
\[
 R_{\Ga_{2m+1}}(\{0,1,2,20\})\le m-1.
\]
All listed centers are distinct legal Fibonacci words in their indicated dimensions.
\end{theorem}
\begin{proof}
Write $w=\mathrm{wt}(v)$ for a legal word $v$. The occupied coordinates
form an independent set of the path on $n$ coordinate positions.

If $n=2m$, then $w\le m$. For $w\le m-1$, the center $0$
is within distance $m-1$. For $w=m$, coordinates $0$ and $1$ cannot
both be zero: the remaining $2m-2$ positions contain at most $m-1$
nonadjacent ones. Thus $v$ has a one in coordinate $0$ or $1$,
and $d(v,1)=m-1$ or $d(v,2)=m-1$.

If $n=2m+1$, then $w\le m+1$. The cases $w\le m-1$ and
$w=m$ with coordinate $0$ or $1$ occupied are covered as before.
If $w=m$ and both first coordinates are zero, positions
$2,\ldots,2m$ form a path of odd length $2m-1$; attaining its
maximum independent-set size $m$ forces exactly coordinates
$2,4,\ldots,2m$ to be occupied. In particular coordinates $2$ and $4$
are both occupied. Since $20=2^2+2^4$,
$d(v,20)=m-2$. Finally $w=m+1$ is possible only for the alternating
word with ones exactly in positions $0,2,\ldots,2m$.
It also contains positions $2,4$, so $d(v,20)=m-1$.
All cases are covered.
\end{proof}

\begin{corollary}[finite-radius covers below the cop number]\label{cor:gap}
For every $m\ge6$, $\Ga_{2m}$ admits a radius-$(m-1)$ covering
with three centers even though $c(\Ga_{2m})>3$.
For every $m\ge7$, $\Ga_{2m+1}$ admits such a covering with four
centers even though $c(\Ga_{2m+1})>4$.
Consequently, these static coverings cannot certify an eventual
winning cop strategy, let alone a capture-time upper bound.
\end{corollary}
\begin{proof}
Theorem~\ref{thm:cover} supplies the covers. Formula~\eqref{eq:cop}
gives $c(\Ga_{2m})=\lceil(m+1)/2\rceil>3$ when $m\ge6$,
and $c(\Ga_{2m+1})=\lceil(m+2)/2\rceil>4$ when $m\ge7$.
\end{proof}

\noindent\textbf{Scope of the corollary.}
This is a comparison between the number of covering centers and the cop
number, not a proved infinite-family inequality between covering radius
and capture time. Indeed, the single center $0$ covers $\Ga_{2m}$
with radius $m$ and $\Ga_{2m+1}$ with radius $m+1$, since their
largest Hamming weights are $m$ and $m+1$, respectively. Theorem
\ref{thm:cover} improves these elementary one-center radii by just
one (even case) or two (odd case). In particular, the corollary does
\emph{not} establish $m-1<\capt(\Ga_n)$ on either infinite family.

\section{Two concrete failed certification routes}
The corollary yields a cop-number comparison only, not an
infinite-family capture-time separation. The following two small
instances explain precisely what failed in No.72.

\begin{proposition}[the dimension-ten separation]\label{prop:ten}
The opening $\{0,1,2\}$ radius-four covers all $144$ vertices
of $\Ga_{10}$, but it cannot guarantee capture within four moves.
Indeed, the inherited independently audited finite result gives
$\capt_3(\Ga_{10})=5$.
\end{proposition}
\begin{proof}
Take $m=5$ in Theorem~\ref{thm:cover}. The second assertion uses
the archived v0.18 complete finite-game certificate and its independent
audit; it does not follow from the covering calculation. The theorem
provides a self-contained proof of the covering part, whereas the
capture-time statement retains its explicit computer-assisted dependency.
\end{proof}

\begin{proposition}[dimension eleven: simultaneous feasibility]\label{prop:eleven}
Let $C=\{0,1,2,20\}\subseteq\Ga_{11}$ and
$H=\{x\in\Ga_{11}:x_i=0\text{ for odd }i\}$.
Then $H\cong Q_6$ is a retract, $c(H)=4$, and
\[
 R_{\Ga_{11}}(C)\le4,\qquad
 \bigl(d(c,H):c\in C\bigr)=(0,0,1,0).
\]
Thus both the four-move radius-covering requirement and the
particular v0.22 delayed-retract necessary condition are simultaneously
satisfied by $C$.
\end{proposition}
\begin{proof}
The covering follows from Theorem~\ref{thm:cover} with $m=5$.
The six even-indexed coordinates are independent, and erasing all odd
coordinates maps legal vertices and edges to legal vertices and edges
(or stays), fixing $H$ pointwise. It is therefore a retraction onto
$Q_6$. The standard hypercube cop-number formula gives $c(Q_6)=4$.
Distances to $H$ count occupied odd coordinates by Lemma~\ref{lem:metric}.
For the four listed centers these counts are $0,0,1,0$.
\end{proof}

The v0.22 delayed-retract rule excludes an opening only when some cop
starts at distance strictly greater than the target turn limit from a
retract whose cop number equals the number of available cops. Since no
cop in Proposition~\ref{prop:eleven} is farther than four from $H$,
that rule excludes none of the four centers. This is a \emph{failure of
the conjunction of these particular necessary-condition tests to prove
nonexistence of all four-turn openings}; it is \emph{not} a proof that
$C$ wins in four moves. A stronger dynamic criterion might still
eliminate $C$ or all openings.

\section{What the failed route does and does not show}
We record the outcome of the final research trial without changing the
game's currently unresolved dimension-eleven status.
\begin{center}
\small
\begin{tabular}{p{0.58\linewidth}p{0.29\linewidth}}
\toprule
Claim or verification target & Status \\
\midrule
Explicit even/odd infinite-family static covers & Proved (elementary) \\
Covering with fewer centers than required cops & Proved using known cop-number formula \\
$\Ga_{10}$ static/dynamic separation & Cover proved; finite game audited in v0.18 \\
Infinite-family radius below $\capt(\Ga_n)$ & \textbf{Not established} \\
$\Ga_{11}$: cover and specific delayed-retract test both pass & Proved \\
$\capt(\Ga_{11})\ge5$ & \textbf{Not proved} \\
All possible structural lower-bound methods fail & \textbf{Not claimed} \\
Originality of the new presentation or formulas & Inconclusive \\
Peer review / novelty assessment & None / inconclusive \\
\bottomrule
\end{tabular}
\end{center}

The inherited bound remains
\[
 4\le\capt(\Ga_{11})\le17,
\]
with the limitations of the inherited lower-bound evidence unchanged.
The withdrawn v0.8 five-move argument is invalid and must not be quoted
as proof. Our present negative result is about \emph{method adequacy}, not
about the exact value of $\capt(\Ga_{11})$.

\section{Reproducibility and novelty assessment}
The accompanying Python verifier constructs each $\Ga_n$ and performs
an independent breadth-first-search distance check for the listed centers
at every dimension $2\le n\le19$ for which the stated construction applies.
It additionally checks the $233$ vertices of $\Ga_{11}$, the $64$-vertex
retract $H$, and the four distances appearing in
Proposition~\ref{prop:eleven}. It checks SHA256 hashes of frozen
prior-study evidence. The general results rely on the written proofs;
finite testing is only a sanity check. The verifier does \emph{not}
rerun a four-cop $\Ga_{11}$ minimax game or independently prove the
archived $\Ga_{10}$ capture-time assertion.

\textbf{Novelty: INCONCLUSIVE.} This note's elementary ball covers and
the logical distinction between distance covering and winning pursuit
may be implicit in established covering-code and pursuit-evasion
literature. The author claims no priority and no new capture-time formula.
References have been checked selectively, not exhaustively.
A future mathematical paper would require a sufficiently distinct
result and a separate novelty review. The appropriate classification of
this draft is a negative research record, not a validated discovery of
a universal no-go theorem.

\vspace{3pt}\noindent\textbf{Research disposition.}
No.72, version v1.0 (negative research record): \textbf{not peer reviewed; novelty inconclusive}.
The static/retract route is discontinued for lack of leverage.
This record is intended to document a failed proof approach, not a new lower bound for $\Ga_{11}$ or an established priority claim.
\par\smallskip
\noindent\textcopyright\ 2026 Ryutaro Yonezu. All rights reserved.

\begin{thebibliography}{9}
\bibitem{hxx2026}
Z. Huang, Y.-T. Xie, and S.-J. Xu,
\emph{Cop numbers for subclasses of partial cubes},
arXiv:2608.25416 (2026), \url{https://arxiv.org/html/2608.25416}.
\bibitem{munarini2002}
E. Munarini and N. Zagaglia Salvi,
\emph{Structural and enumerative properties of the Fibonacci cubes},
Discrete Mathematics 255 (2002), 317--324,
\url{https://doi.org/10.1016/S0012-365X(01)00407-1}.
\bibitem{azarija2018}
J. Azarija, S. Klav\v{z}ar, Y. Rho, and S. Sim,
\emph{On domination-type invariants of Fibonacci cubes and hypercubes},
Ars Mathematica Contemporanea 14 (2018), 387--395,
\url{https://doi.org/10.26493/1855-3974.1172.bae}.
\bibitem{bgkp2013}
A. Bonato, P. Gordinowicz, W. B. Kinnersley, and P. Pra\l at,
\emph{The capture time of the hypercube},
Electronic Journal of Combinatorics 20(2), P24 (2013),
\url{https://doi.org/10.37236/2921}.
\bibitem{bonato2009}
A. Bonato, G. Hahn, P. A. Golovach, and J. Kratochv\'il,
\emph{The capture time of a graph},
Discrete Mathematics 309 (2009), 5588--5595,
\url{https://doi.org/10.1016/j.disc.2008.04.004}.
\bibitem{over2016}
A. Bonato et al., \emph{The game of overprescribed Cops and Robbers played on graphs},
arXiv:1611.07592 (2016), \url{https://arxiv.org/abs/1611.07592}.
\end{thebibliography}
\end{document}
